\section{Biology-Informed Failure Modes}
\label{sec:failure_modes}

This is the analytical core of the paper. We document four failure modes encountered during development of MI. For each, we describe (i) the symptom, (ii) the biological parallel, (iii) the fix derived from biology, and (iv) the empirical impact of the fix.

\subsection{Modulatory Explosion}

\paragraph{Symptom.}
A positive feedback loop in which high reward triggers strong dopamine release, which potentiates active synapses, which raises future activity, which increases reward, which triggers more dopamine. All motor morphons saturate at maximum firing rate within $\sim$50 ticks; the reward signal becomes meaningless.

\paragraph{Biological parallel.}
Dopaminergic excitotoxicity. Biology prevents this through receptor desensitization on sustained exposure, presynaptic autoreceptor inhibition, and synaptic reuptake.

\paragraph{Fix.}
Receptor saturation via Hill function ($n=2$), receptor downregulation on sustained exposure, exponential decay (``reuptake'') of modulator concentrations from release sites.

\paragraph{Result.}
Stable modulatory dynamics. Dopamine signal retains information content; learning continues without saturation.

\subsection{Motor Silencing from Burst Dead Zone}

\paragraph{Symptom.}
After an initial burst phase, motor morphons stop firing entirely. Network becomes unresponsive even when sensory input is strong. Forward pass produces zero output regardless of state.

\paragraph{Biological parallel.}
Depolarization block. In biology, intrinsic excitability homeostasis (slow conductance modulation that keeps neurons in their dynamic range) and tonic baseline activity prevent complete silencing \cite{turrigiano1998activity}.

\paragraph{Fix.}
Tonic baseline current injection into motor morphons (an analogue of intrinsic excitability), and Endoquilibrium's threshold-bias rule which lowers firing thresholds when associative firing rate drops below its setpoint.

\paragraph{Result.}
Motor morphons maintain baseline activity. STDP remains effective across firing rate regimes. The CartPole solve depends on this fix; without it, the system stalls at $\sim$10 steps (random baseline) within 50 episodes.

\subsection{LTD Vicious Cycle}

\paragraph{Symptom.}
When the LTD component of STDP is too strong relative to LTP, weakened synapses carry less current, which reduces post-synaptic firing, which increases the LTD/LTP ratio (because post-before-pre events dominate at low firing rates), which weakens synapses further. Global weight decay to near-zero. Network falls silent. All morphons die from energy starvation.

\paragraph{Biological parallel.}
Runaway LTD / synaptic scaling failure. In biology, homeostatic synaptic scaling \cite{turrigiano2008homeostatic} and BCM-like metaplasticity thresholds \cite{bienenstock1982bcm} prevent this.

\paragraph{Fix.}
Synaptic scaling: every $\sim$50 ticks, normalize the total incoming weight to each morphon to a target. This is the Turrigiano operation, implemented in the homeostasis subsystem. BCM-style metaplasticity thresholds adjust the LTP/LTD crossover point based on recent post-synaptic activity, preventing the runaway LTD.

\paragraph{Result.}
Weight distribution stays bounded. No more global silencing. Networks survive multi-hour training runs without entering the death spiral.

\subsection{Premature Mature: A Novel Failure Mode}
\label{sec:premature_mature}

The fourth failure mode is one we have not seen documented in prior work, and it is broadly applicable to any biologically-inspired homeostatic regulator that uses raw reward as a maturity signal.

\paragraph{Symptom.}
On MNIST, accuracy plateaus at $\sim$26\% regardless of training duration. The associative layer has formed digit-selective firing patterns (verified by post-hoc neuron labeling), but the supervised readout cannot exploit them. Plasticity has been throttled to $\text{plasticity\_mult}=0.60$ within the first 100 training samples. Endoquilibrium reports stage = \emph{Mature}.

\paragraph{Diagnosis.}
The Endoquilibrium stage detector triggers \emph{Mature} when the slow EMA of reward exceeds 0.05 with low variance. This logic was calibrated for reinforcement learning tasks (CartPole, where reward is sparse and correlated with actual performance: reward only reaches high values when the system is genuinely competent). It fails on supervised classification for two reasons:

\begin{enumerate}
  \item \textbf{Dense reward injection.} The supervised training loop calls \texttt{reward\_contrastive(0.2)} on \emph{every} sample regardless of whether the prediction was correct. Reward is no longer a signal of mastery --- it is a constant. The slow EMA converges to $\sim$0.65 within 500 samples \emph{regardless of accuracy}, far above the 0.05 threshold.
  \item \textbf{Inflated startup window.} The performance signal Endo tracks (\texttt{recent\_performance}) is fed by a windowed accuracy. During the first $\sim$50 samples the window has not filled, so a single correct guess gives accuracy 1.0, then 0.5, then 0.33$\ldots$. These early high values get integrated by the $\alpha = 0.05$ EMA before the window settles to the true accuracy ($\sim$0.10), pumping the EMA above the Mature threshold within the warmup period.
\end{enumerate}

Both pathways converge to the same outcome: the slow EMA reaches 0.65 by sample 500, low variance, low trend $\to$ Endo declares Mature $\to$ \texttt{plasticity\_mult} drops to 0.60 $\to$ the supervised gradient is multiplied by 0.6 at every weight update $\to$ learning effectively stops at $\sim$26\%. The hidden layer has formed digit-selective firing patterns (verified by post-hoc neuron labeling), but the readout cannot exploit them because the gradient that would tune it is being throttled by a regulator that thinks the system is done learning.

\paragraph{Empirical evidence.}
\begin{table}[h]
\centering
\caption{MNIST accuracy as a function of plasticity\_mult (pm). Each row is a separate run with the same network and seed.}
\label{tab:plasticity}
\begin{tabular}{lcc}
\toprule
Plasticity level & Endo stage & Accuracy \\
\midrule
pm = 0.60 (Mature, premature)        & Mature                          & 27\% \\
pm = 0.96--1.77 (oscillating)        & Differentiating/Consolidating  & \textbf{31\%} \\
pm = 1.80 (Differentiating, constant) & Differentiating                & 25\% \\
pm = 2.16 (post-damage recovery)      & Differentiating                & \textbf{52.5\%} \\
\bottomrule
\end{tabular}
\end{table}

The pattern is monotonic in plasticity \emph{within the oscillating regime} (Figure~\ref{fig:plasticity_accuracy}): more plasticity, more accuracy. But pure constant high plasticity (1.80) also fails (25\%) because the weights churn without ever consolidating. The system needs both phases.

\begin{figure}[h]
  \centering
  \figifexists[0.7\linewidth]{plasticity_accuracy.pdf}{plasticity vs accuracy scatter}
  \caption{MNIST accuracy as a function of plasticity multiplier and Endoquilibrium stage. Premature Mature (pm=0.60) caps accuracy at 27\%. The natural Differentiating/Consolidating oscillation around pm=1.4 reaches 31\%. Constant high plasticity (pm=1.80) is \emph{worse} (25\%) because weights churn without consolidating. Post-damage recovery at pm=2.16 reaches 52.5\%. The relationship is monotonic only when the system can also consolidate.}
  \label{fig:plasticity_accuracy}
\end{figure}

\paragraph{Biological parallel.}
This is reward prediction error vs.\ raw reward as a maturity signal. In biology, the brain does not declare ``learning complete'' because rewards are arriving --- it tracks reward \emph{prediction error} (RPE), which goes to zero only when actual outcomes match expected outcomes. A constant 0.2 reward is only ``mature'' if the system was \emph{predicting} 0.2; if the system is at 11\% accuracy but receiving 0.2 each sample, RPE is high and learning should continue. The biological literature on dopamine and RPE \cite{schultz1997neural} is unambiguous on this point.

\paragraph{Fix.}
Two changes:
\begin{enumerate}
  \item Gate the Mature transition on a monotonically-increasing tick counter (\texttt{total\_updates}) rather than the deduplicating reward history buffer (which is capped at 500 entries and so could never satisfy a 2000-tick gate). The new gate requires \texttt{total\_updates} $\geq$ 2000 before Mature can trigger. This prevents Mature from firing during the critical learning phase regardless of how inflated the reward signal becomes. Empirically, this prevents the throttle and lets the system learn through the natural Differentiating$\leftrightarrow$Consolidating oscillation.
  \item (Architectural, future work.) Replace the raw reward signal feeding Endo with reward prediction error from the planned Limbic Circuit's Motivational Drive component. This is the principled fix; the tick gate is a patch that buys time until RPE-based regulation is implemented.
\end{enumerate}

\paragraph{The oscillation discovery.}
A surprising secondary finding: pure high plasticity is worse than oscillating plasticity. When we tested raising both the Mature \emph{and} Consolidating gates (preventing all throttling), accuracy dropped to 25\%. The system needs the natural Differentiating$\leftrightarrow$Consolidating oscillation that biology calls the theta/sharp-wave-ripple alternation \cite{buzsaki2015hippocampal}: high-plasticity exploration phases interleaved with consolidation phases that lock in gains. The oscillation \emph{is} the feature; the bug was Mature triggering prematurely and locking the system out of the oscillation entirely.

\paragraph{Implications beyond MI.}
Any biologically-inspired homeostatic system that uses raw reward (not RPE) as a maturity signal will exhibit this failure on dense-reward tasks. We expect this to be a common bug in spiking neural network training pipelines that combine RL-style modulation with supervised classification. The fix is to either (i) use sparse reward correlated with correctness, (ii) add a long history gate, or (iii) use RPE.

\subsection{Diagnostic Instruments}
\label{sec:diagnostics}

Each of the four failure modes above was discovered through a specific measurement, not by reading code or guessing. We document the diagnostic instruments here because they are reusable across spiking architectures.

\begin{description}
  \item[\textbf{Jaccard similarity of firing patterns.}] Present the same input twice with voltage reset; the Jaccard similarity of the two fired-morphon sets measures determinism. Values $\geq 0.9$ are stable, $< 0.7$ predict that no readout will converge. Activity stabilization fixes (reducing per-step noise from $\pm 0.05$ to $\pm 0.02$) brought CartPole from Jaccard = 0.43 to 1.0 and were a prerequisite for the SOLVED result.
  \item[\textbf{Output discrimination $d_\pm$.}] $d_+ = V_1 - V_0$ measured for a right-leaning pole probe; $d_-$ for left-leaning. When $d_+$ and $d_-$ have \emph{opposite} signs and $|d| > 5$, the readout is discriminative; when same-signed, the policy is collapsed. This metric surfaced the readout architecture failure (Section~\ref{sec:implementation}) faster than any other diagnostic.
  \item[\textbf{Endoquilibrium summary line.}] The seven channel values (\texttt{rg, ng, ag, hg, tb, pm, hp}) plus the detected stage are logged each medium tick. \texttt{pm} (plasticity multiplier) at 0.6 with stage = Mature on a fresh task is the signature of premature throttling.
  \item[\textbf{Weight standard deviation.}] Synapse weight std $\sigma_w$ over the entire network. $\sigma_w < 0.15$ is feature collapse (all morphons look the same); $\sigma_w > 1.0$ is healthy diversity. Surfaces the LTD vicious cycle and the per-tick capture saturation problem.
\end{description}

\paragraph{Negative results worth noting.}
Several promising-sounding interventions did not work and we document them so others do not retry them: (i) CMA-ES over 15 hyperparameters across 560 evaluations could not exceed avg=21 on CartPole, indicating the bottleneck is architectural rather than parametric; (ii) extending CartPole training from 1000 to 3000 episodes did not improve the asymptote, ruling out ``needs more time''; (iii) TD-scaled learning rate (LR proportional to $|\delta|$) reaches avg=33 versus constant LR avg=195, because TD error is near zero at steady state where the policy is actually learned; (iv) freezing the MI network and training only the readout fails because uniform potentials provide no discriminative signal --- the MI network's diversity is necessary; (v) weight normalization clamping to $[0.5, 2.0]$ collapses weights to uniformity; the gentler $[0.9, 1.1]$ preserves diversity. The pattern across all five negative results is that the readout and the hidden layer are genuinely co-dependent: neither alone is sufficient, neither can be trivially replaced.
